\documentclass[11pt,reqno]{amsart}
\usepackage[margin=1in]{geometry}
\usepackage[expansion=false]{microtype}
\usepackage{xcolor}
\usepackage{amssymb,mathtools}
\usepackage[colorlinks=true,linkcolor=blue!50!black,citecolor=blue!50!black,urlcolor=blue!50!black]{hyperref}
\newtheorem{theorem}{Theorem}[section]
\newtheorem{proposition}[theorem]{Proposition}
\newtheorem{lemma}[theorem]{Lemma}
\newtheorem{corollary}[theorem]{Corollary}
\theoremstyle{remark}
\newtheorem{remark}[theorem]{Remark}
\newcommand{\R}{\mathbb R}
\newcommand{\E}{\mathbb E}
\newcommand{\one}{\mathbf 1}
\newcommand{\ip}[2]{\langle #1,#2\rangle}
\newcommand{\norm}[1]{\lVert #1\rVert}
\newcommand{\ga}{\gamma}
\newcommand{\LL}{\mathcal L}
\newcommand{\HH}{\mathcal H}
\newcommand{\DD}{\mathcal D}
\newcommand{\VV}{\mathcal V}
\newcommand{\RR}{\mathcal R}
\DeclareMathOperator{\tr}{tr}
\DeclareMathOperator{\dist}{dist}
\hypersetup{pdftitle={Plurality Is Stablest for Small Positive Correlations: A Gaussian Simplex Theorem},pdfauthor={Roy van Rijn}}
\title[Gaussian simplex stability at small correlation]{Plurality Is Stablest for Small Positive Correlations:\\A Gaussian Simplex Theorem}
\author{Roy van Rijn}
\subjclass[2020]{60G15, 49K21, 28A75}
\keywords{Gaussian noise stability, regular simplex, plurality is stablest, first moments}
\date{}
\begin{document}
\begin{abstract}
For every fixed number of cells, we prove that the central regular-simplex
partition maximizes Gaussian noise stability among all equal-measure
measurable partitions, for every sufficiently small positive correlation.
The correlation interval is independent of the ambient dimension, and
equality determines the simplex partition up to orthogonal maps, labels
and cylindrical extension. The proof combines Zhang's first-moment
comparison with a compactness argument for arbitrary changes of labels.
We normalize those changes by their exact Ornstein--Uhlenbeck score loss.
A weighted Gaussian estimate prevents this normalized mass from escaping
to infinity; its limiting facet flux has a sharp quadratic cost and can
only be rotational. Alignment of the cell moments removes that mode.
The threshold may depend on the number of cells and is not explicit.
\end{abstract}
\maketitle

\section{The theorem and its mechanism}

Two correlated Gaussian observations need not fall in the same cell of a
partition. Which equal-measure partition makes agreement most likely?
The Gaussian Standard Simplex Conjecture predicts that the answer is the
partition into the winning regions of a centered regular simplex. We prove
this prediction for sufficiently small positive correlation, for each
fixed number of cells.

Let $\ga_d$ be standard Gaussian probability on $\R^d$, with density
$\phi_d$. For $0<\rho<1$, define
\[
 T_\rho f(x)=\int f(\rho x+\sqrt{1-\rho^2}\,y)\,d\ga_d(y),\qquad
 S_\rho(A)=\sum_{i=1}^k\int 1_{A_i}T_\rho1_{A_i}\,d\ga_d.
\]
A measurable partition $A=(A_1,\ldots,A_k)$ is \emph{balanced} if
$\ga_d(A_i)=1/k$ for every $i$. Put $n=k-1$. For vertices
$v_1,\ldots,v_k$ of a centered regular simplex in $\R^n$, let
\[
 T_i=\{x:v_i\cdot x\ge v_j\cdot x\text{ for every }j\}.
\]
In higher dimension $T$ denotes the cylindrical extension of this
partition. All partitions and equality statements are understood modulo
Gaussian null sets.

\begin{theorem}\label{thm:main}
For every integer $k\ge2$ there exists $\rho_k>0$, independent of
$d\ge k-1$, such that every balanced measurable $k$-cell partition of
$\R^d$ satisfies
\[
 S_\rho(A)\le S_\rho(T)\qquad(0<\rho\le\rho_k).
\]
Equality holds precisely for orthogonal images and relabelings of the
cylindrical central regular-simplex partition.
\end{theorem}

The same upper bound holds in dimensions $d<k-1$, by cylindrical
embedding, and is then strict: equality would require a cell-moment
matrix of rank $k-1$. No interval uniform over all $k$ is asserted.

The question is the continuous Gaussian problem associated with Plurality
is Stablest. The two-cell case is classical; see \cite[Introduction]{HT}.
Heilman \cite{Heilman14} proved the three-cell small-correlation theorem
in fixed dimension. Heilman--Tarter \cite[Theorem~1.9]{HT} showed that
every maximizing $k$-cell partition in dimension $d\ge k-1$ is
cylindrical over $\R^{k-1}$, making the three-cell interval dimension
independent. Their theorem reduces maximizers, not individual competitors.
For four cells, the first-moment obstruction to a perturbative extension
was identified in \cite[Section~8, Conjecture~3]{Heilman14}.

Zhang's equal-cell first-moment theorem and covariance comparison
\cite{Zhang} supply that endpoint for every $k$. We give the covariance
argument, with attribution, in Section~\ref{sec:alignment}; its geometric
input is the Gaussian multi-bubble theorem of Milman--Neeman
\cite[Theorem~1.1]{MN}. The additional task is nonlinear: a partition
close in Gaussian measure can still move arbitrarily small amounts of
mass arbitrarily far away. A second variation for smooth interfaces does
not by itself compare such partitions.

\subsection*{The right scale for a change of labels}
Write $e_i=1_{A_i}-1_{T_i}$. The exact comparison is
\begin{equation}\label{eq:comparison-intro}
 S_\rho(A)-S_\rho(T)=\widetilde Q_\rho(e)-2\LL_\rho(A),
\end{equation}
where $\LL_\rho(A)$ is the loss incurred by assigning the labels of
$A$ to the fixed scores $T_\rho1_{T_i}$, and $\widetilde Q_\rho$ is
the quadratic interaction of the changes, with constants removed.
Both terms matter. The first-moment deficit localizes a hypothetical
competitor, but cannot serve as a uniform scale for its full noise loss;
Proposition~\ref{prop:obstruction} gives a remote-ball example.

The local picture explains our normalization. Let $t$ be distance into
an old cell across one of its facets, and let $\ell$ be the common
distance between adjacent cell moments. For small correlation, the
score difference is $\rho\ell t$ to first order. Moving a one-sided
strip of width $h$ therefore costs $\rho\ell h^2/2$ per unit Gaussian
surface area. Its normalized change of measure is a facet flux of
density $h$. This elementary strip calculation predicts the quadratic form
\[
 Q(f)=\ell\norm{f}_{\HH}^2-\norm{\DD f}_F^2:
\]
the first term is the cost of moving mass and the second its first-moment
interaction. Quantitative moment alignment shows that $Q$ is nonnegative
and its only zero modes are rotations.

To pass from strips to arbitrary measurable changes, we normalize by
$s^2=\LL_\rho(A)/\rho$. A weighted Gaussian estimate keeps the
normalized transfer measures tight, even if $s$ is smaller than every
power of $\rho$. Their limiting facet flux satisfies
$\norm{\DD f}_F^2\ge2\ge\ell\norm{f}_{\HH}^2$ and must therefore
be rotational. Choosing the closest regular cell-moment matrix makes
its moment change exactly orthogonal to rotations. This yields the
contradiction.

Switching-surface curvature is an established second-order method;
Christof--Wachsmuth \cite[Theorems~4.4 and~6.12]{CW} give general
conditions for local quadratic growth. Their explicit bang-bang theorem
assumes a bounded domain, a scalar constraint and a regular switching
surface. The Gaussian multicell setting also has unbounded tails and
rotational degeneracy. We prove the needed compactness and flux estimates
directly in Sections~\ref{sec:tails}--\ref{sec:compactness}.

\section{First moments and quantitative alignment}\label{sec:alignment}

Fix $k$ and put $n=k-1$. Constants below may depend on $k$, but not on
an ambient dimension $d\ge n$. Let
\[
 P=I-\one\one^T/k,\quad E=\one^\perp\subset\R^k,\quad Q_*=P/n,
 \qquad F=\frac{(\E\max_{i\le k}Z_i)^2}{n},
\]
where the $Z_i$ are independent standard normals. For a partition let
$B$ have rows $b_i=\int_{A_i}x\,d\ga_d$; then
$W_1(A)=\norm B_F^2$ and $\sum_i b_i=0$.
The regular moment orbit is
$\mathcal M_d=\{B_*:B_*B_*^T=FQ_*\}$. Every member is the moment
matrix of its central regular winning partition.

For centered positive semidefinite $Q$, define the balanced score value
\[
 \mathcal C(Q)=\inf_{\lambda\in\R^k}
 \left\{k^{-1}\sum_i\lambda_i+\E\max_i(Y_i-\lambda_i)\right\},
 \qquad Y\sim N(0,Q).
\]
Here centered means $Q\one=0$. When $Q$ is positive on $E$, minimizing
prices make the affine winning cells have masses $1/k$. Write
$\ell_{ij}$ for the score-site distance, $A_{ij}$ for the ordinary
Gaussian facet area, and $w_{ij}=A_{ij}/\ell_{ij}$.

\begin{proposition}[Zhang \cite{Zhang}]\label{prop:zhang}
Every balanced measurable Gaussian partition satisfies $W_1(A)\le F$.
Equality holds precisely for central regular-simplex partitions and their
cylindrical extensions; it is impossible below dimension $n$.
For every centered trace-one $Q\succeq0$, set $Q_t=Q_*+t(Q-Q_*)$.
Then
\begin{equation}\label{eq:cov-deficit}
 F-\mathcal C(Q)^2=\int_0^1\frac{D(t)}{t^2}\,dt,\qquad
 D(t)=\frac2n\left[\sum_{e<f}w_e(t)w_f(t)(\ell_e(t)-\ell_f(t))^2
                         +S(t)^2-S_*^2\right],
\end{equation}
where $e,f$ range over distinct unordered label pairs,
$S(t)=\sum_{i<j}A_{ij}(t)\ge S_*=\sqrt{nF/2}$, and all weights
are positive for $t<1$.
\end{proposition}

\begin{proof}
We follow \cite[Sections~2--5]{Zhang}. Two facts include singular
covariances: $\mathcal C(sQ)=\sqrt s\,\mathcal C(Q)$, and
$\mathcal C$ is continuous. For the latter, couple $Q^{1/2}G$ and
$\widehat Q^{1/2}G$; the expected norm of their difference bounds the
difference of the two price objectives uniformly in the prices.

For full rank on $E$, the score-difference map is nonsingular. Balancing
prices exist: in the gauge $\sum_i\lambda_i=0$, the objective is at
least $-\min_i\lambda_i$ and is coercive. Its gradient gives the mass
errors. Each pair has a positive facet, and the price Hessian is the
weighted Laplacian $L$, with edge weights $w_{ij}$. Thus prices are
unique and smooth in the zero-sum gauge. The same score coordinates
turn the facet domains into fixed polyhedra; Gaussian domination gives
continuity of the weights and smoothness of the score integrals.

For a score matrix $M$ with $MM^T=Q$, Gaussian flux gives its winning
moment matrix $LM$. Truncation at large balls justifies flux on these
unbounded cells. The envelope derivative and this identity give
\begin{equation}\label{eq:cov-derivative}
 \mathcal C(Q)=\tr(LQ)=\sum_{i<j}w_{ij}\ell_{ij}^2,
 \qquad d\mathcal C(Q)[H]=\tfrac12\tr(LH).
\end{equation}
Indeed write $H=MAM^T$, $A=A^T$, and use
$M(t)=M(I+tA)^{1/2}$; its derivative is $MA/2$.

For $Q_*$, regular scores have distribution
$(Z_i-\overline Z)/\sqrt n$, so $\mathcal C(Q_*)=\sqrt F$.
Symmetry in \eqref{eq:cov-derivative} gives $L_*=\sqrt F P$ and
$S_*=\sqrt{nF/2}$. The Gaussian multi-bubble theorem
\cite[Theorem~1.1]{MN} gives $S\ge S_*$ for the balanced cells in
dimension $n$; its cell-count hypothesis is exactly $k=n+1$.
Perimeter here is the sum of interface areas, half the sum of individual
cell perimeters. Weighted Cauchy--Schwarz now yields
\[
 \mathcal C(Q)\tr L\ge2S^2\ge nF.
\]
Along $Q_t$, write $C(t)=\mathcal C(Q_t)$ and $h(t)=C(t)^2-F$.
Equation~\eqref{eq:cov-derivative} gives
\[
 tC'(t)=\tfrac12(C(t)-\tr L_t/n),\qquad
 (h(t)/t)'=-D(t)/t^2\le0.
\]
Expanding the weighted Cauchy difference gives exactly $D(t)$ in
\eqref{eq:cov-deficit}. Since $L_*=\sqrt F P$ and
$\tr(Q-Q_*)=0$, we have $C'(0)=0$ and $h(t)/t\to0$ at zero.
Integration and continuity at $t=1$ prove the identity, including its
improper endpoint, and $\mathcal C(Q)\le\sqrt F$.
Equality forces all pair distances equal at an interior time, hence
$Q=Q_*$. Centered trace-one equal-distance scores have Gram matrix
$P/n$; for $k=2$ the covariance domain is already this singleton.

Finally the original labels give a feasible balanced assignment:
$W_1(A)\le\mathcal C(BB^T)$. If $W_1(A)>0$, homogeneity gives
$\sqrt{W_1(A)}\le\mathcal C(BB^T/W_1(A))\le\sqrt F$.
Equality forces regular score covariance and assignment equality.
Regular balancing prices are equal by symmetry, ties are null, and the
labels must be the regular winners. Their actual moment matrix follows
from the flux identity above. Rank excludes equality when $d<n$.
\end{proof}

\begin{proposition}\label{prop:alignment}
There are $L,K<\infty$ such that, for every balanced partition in
$d\ge n$, a closest $B_*\in\mathcal M_d$ satisfies
\begin{equation}\label{eq:alignment}
 \norm{B-B_*}_F^2\le L\delta,\qquad
 \sum_i\ga_d(A_i\triangle T_i)\le K\sqrt\delta,
 \qquad \delta=F-W_1(A).
\end{equation}
Here $T$ is the winning partition of $b_i^*\cdot x$. The same choice
has $B-B_*\perp T_{B_*}\mathcal M_d$.
\end{proposition}

\begin{proof}
We first obtain covariance coercivity. For $0\le t\le1/2$, every
eigenvalue of $Q_t$ on $E$ lies in $[1/(2n),1]$. This is a compact
set inside the positive cone. The continuous positive facet weights
therefore have a common lower bound $\alpha>0$.
Put $H=Q-Q_*$, $d_{ij}=(e_i-e_j)^TH(e_i-e_j)$ and $q=k(k-1)/2$.
Then $\ell_{ij}(t)^2=2/n+td_{ij}$, $\ell_{ij}(t)\le\sqrt2$, and
$\sum_{i<j}d_{ij}=0$. If $D_H$ has these off-diagonal entries and
zero diagonal, $H=-PD_HP/2$. Consequently
\[
 \sum_{e<f}(d_e-d_f)^2=q\sum_e d_e^2\ge2q\norm H_F^2,
 \qquad (\ell_e(t)-\ell_f(t))^2\ge t^2(d_e-d_f)^2/8.
\]
Keeping $0<t<1/2$ in \eqref{eq:cov-deficit} proves
\begin{equation}\label{eq:cov-coercivity}
 F-\mathcal C(Q)^2\ge\beta\norm{Q-Q_*}_F^2,
 \qquad \beta=q\alpha^2/(4n)>0.
\end{equation}
For $k=2$ the right side is zero, so the same conclusion holds.

Write $W=\norm B_F^2$. Feasible assignment and homogeneity give
$\norm{BB^T/W-Q_*}_F^2\le\delta/\beta$ when $W>0$.
Choose the closest regular matrix by singular value decomposition.
For the $n$ singular values $s_j$ of $B$ and $a=\sqrt{F/n}$,
\[
 \norm{B-B_*}_F^2=\sum_j(s_j-a)^2
 \le\frac nF\norm{BB^T-FQ_*}_F^2
 \le(1+nF/\beta)\delta.
\]
Here the covariance difference splits orthogonally into its trace and
trace-zero parts. Complete right singular vectors if $B$ has deficient
rank; if $W=0$, the squared distance is simply $F=\delta$.

The regret against these same regular scores is exactly
\begin{equation}\label{eq:linear-regret}
 r=\int(\max_i b_i^*\cdot x-b_{I_A(x)}^*\cdot x)d\ga_d
   =F-\ip{B_*}{B}_F
   =\tfrac12(\delta+\norm{B-B_*}_F^2).
\end{equation}
Let $g$ be the gap between the largest and second-largest regular scores.
Each pairwise difference is normal with standard deviation
$\ell=\sqrt{2F/n}$, so a union bound gives $\ga_d(g\le u)\le Mu$
for a constant $M$ depending only on $k$. If
$\eta=\ga_d(I_A\ne I_T)$, then $\eta\le Mu+r/u$.
Minimizing in $u$ and using total symmetric difference $2\eta$ proves
\eqref{eq:alignment}. Orthogonality follows by differentiating squared
distance along the regular moment orbit at its minimum.
\end{proof}

\section{The interface form and its rotational kernel}\label{sec:interface}

We now work in $\R^n$, with $k$ fixed. Let
$n_{ij}=(b_i^*-b_j^*)/\ell$, $H_{ij}=n_{ij}^\perp$,
$F_{ij}=T_i\cap T_j$, and $d\nu_{ij}=\phi_n d\mathcal H^{n-1}$
on the pair plane. Use the same measure on its facet. Let $f_{ij}$ be
outward speed of $T_i$, with velocity $-f_{ij}n_{ij}$, and put
$f_{ji}=-f_{ij}$. Define
\[
 \HH=\bigoplus_{i<j}L^2(F_{ij},\nu_{ij}),\qquad
 \HH_0=\{f:\sum_{j\ne i}\int f_{ij}d\nu_{ij}=0\text{ for every }i\},
 \qquad (\DD f)_i=\sum_{j\ne i}\int x f_{ij}d\nu_{ij}.
\]
The map $\DD$ is bounded and finite rank, by Gaussian facet second
moments. For skew $\Omega$ write
\[
 f^\Omega_{ij}=-(\Omega x)\cdot n_{ij},\qquad
 \RR=\{f^\Omega\},\qquad
 \VV=\{\text{matrices with rows }\Omega b_i^*\}.
\]
In dimension one, pair planes are points with their zero-dimensional
Gaussian surface measure, and $\HH_0=\{0\}$.

\begin{lemma}\label{lem:interface}
On $\HH_0$, $Q(f)=\ell\norm f_\HH^2-\norm{\DD f}_F^2$ is
nonnegative and its kernel is exactly $\RR$, of dimension $n(n-1)/2$.
It has a positive gap modulo that kernel.
\end{lemma}

\begin{proof}
The one-dimensional assertion is immediate. Otherwise begin with smooth
speeds compactly supported in facet interiors and satisfying the linear
mass conditions. Their interface motions admit exact balance: correct
the $O(t^2)$ mass errors using $n$ fixed facet bumps whose label pairs
form a spanning tree. The mass derivative has rank $n$, so the implicit
function theorem gives corrections with zero initial speed.
Boundary differentiation yields
\begin{equation}\label{eq:moment-variation}
 B'(0)=\DD f,\qquad -\tfrac12 W_1''(0)=Q(f).
\end{equation}
Indeed $(b_i^*-b_j^*)\cdot x$ vanishes on the facet and its derivative
under velocity $-f_{ij}n_{ij}$ is $-\ell f_{ij}$. Acceleration terms
are multiplied by the vanishing score, while differentiation of the
squared moments contributes $2\norm{\DD f}_F^2$.
Proposition~\ref{prop:zhang} now gives $Q\ge0$.

For this deformation $B(t)=B_*+t\DD f+O(t^2)$ and
$F-W_1(A(t))=t^2Q(f)+o(t^2)$. Apply Proposition~\ref{prop:alignment}
and take the first-order distance to the smooth moment orbit:
\begin{equation}\label{eq:tangent-alignment}
 \dist_F(\DD f,\VV)^2\le LQ(f).
\end{equation}
These smooth balanced speeds are dense in $\HH_0$: truncate and smooth
on facet interiors, then correct their converging linear mass errors by
the same fixed bumps. Thus both conclusions extend by continuity.

Rotations preserve masses and $W_1$, so $Q(f^\Omega)=0$ and
$\DD f^\Omega$ has rows $\Omega b_i^*$. Gaussian integrability
justifies these identities by truncation. If $Q(f)=0$,
\eqref{eq:tangent-alignment} supplies $u\in\RR$ with $\DD u=\DD f$.
A nonnegative quadratic form has zero mixed pairing with its kernel;
hence $Q(f-u)=0$. Since $\DD(f-u)=0$, this forces $f=u$.
The moment vectors span $\R^n$, proving the dimension claim. Finally
the operator of $Q$ on $\HH_0$ is $\ell I$ minus a finite-rank
operator, so its nonzero spectrum is bounded away from zero.
\end{proof}

\section{Exact score loss and weighted transfer bounds}\label{sec:tails}

Put $\sigma=\sqrt{1-\rho^2}$ and
$U_{\rho,ij}=T_\rho(1_{T_i}-1_{T_j})$. In coordinates
$x=z+tn_{ij}$, $z\in H_{ij}$, let $\phi_0=\phi_1(0)$; then
\[
 d\ga_n=\phi_1(t)dt\,d\ga_{n-1}(z),\qquad
 d\nu_{ij}=\phi_0d\ga_{n-1}(z).
\]
For $n=1$, $\ga_0$ is unit mass on the pair plane.

\begin{lemma}\label{lem:margin}
There is $c_T>0$ such that for every $0<\rho<1$ and $x\in T_i$,
\begin{equation}\label{eq:saturated-margin}
 U_{\rho,ij}(x)\ge c_T\min\{\rho n_{ij}\cdot x/\sigma,1\}.
\end{equation}
Also $U_{\rho,ij}=0$ on $H_{ij}$ and
$\nabla U_{\rho,ij}/\rho\to b_i^*-b_j^*$ locally uniformly as
$\rho\downarrow0$.
\end{lemma}

\begin{proof}
The transposition of vertices $i,j$ is reflection $G$ in $H_{ij}$.
Cone scaling, with $m=\rho x/\sigma$, gives
\[
 U_{\rho,ij}(x)=\int_{T_i}[\phi_n(y-m)-\phi_n(Gy-m)]dy,
 \qquad \frac{\phi_n(Gy-m)}{\phi_n(y-m)}=e^{-2m_Ny_N}.
\]
For $m\in T_i$ the integrand is nonnegative. Fix a ball $D$ in the
interior of $T_i$ with $z_N\ge a>0$ on $D$ and $\ga_n(D)=p>0$.
Since $m+D\subset T_i$, restriction to this ball gives
$p(1-e^{-2am_N})\ge p(1-e^{-2a})\min\{m_N,1\}$.
Choose constants common to the finitely many pairs. Reflection also
gives exact vanishing on the entire pair plane.

The Mehler kernel relative to $\ga_n\otimes\ga_n$ is
\begin{equation}\label{eq:mehler}
 M_\rho(x,y)=(1-\rho^2)^{-n/2}
 \exp\!\left(\frac{2\rho x\cdot y-\rho^2(|x|^2+|y|^2)}{2(1-\rho^2)}\right).
\end{equation}
Its value and correlation derivative at zero are $1$ and $x\cdot y$.
Its first spatial derivatives are Gaussian dominated for $x$ in compact
sets and small $\rho$. Taylor expansion and the cell moments give the
claimed gradient limit.
\end{proof}

For any partition define the ordered transfers
$h_{ij}=1_{T_i\cap A_j}$, $i\ne j$, and
$e_i=\sum_{j\ne i}(h_{ji}-h_{ij})=1_{A_i}-1_{T_i}$. Its exact score
loss is
\begin{equation}\label{eq:regret-cost}
 \LL_\rho(A)=\sum_{i\ne j}\int U_{\rho,ij}h_{ij}d\ga_n
 \ge c_T\rho\sum_{i\ne j}\int\min\{n_{ij}\cdot x,1\}h_{ij}d\ga_n,
\end{equation}
because $\min\{\rho t/\sigma,1\}\ge\rho\min\{t,1\}$.
The loss is positive whenever $A\ne T$ modulo null sets.

\begin{lemma}\label{lem:weighted}
Let $w(x)=(1+|x|^2)e^{|x|^2/16}$. Uniformly over unit vectors $u$
and measurable $0\le h\le1$ supported in $t=u\cdot x\ge0$,
\begin{equation}\label{eq:weighted-tail}
 \left(\int_{|x|>R}wh\,d\ga_n\right)^2
 \le\beta(R)\int\min\{t,1\}h\,d\ga_n\qquad(R\ge2),
\end{equation}
where $\beta(R)\to0$. There is a finite constant for the same estimate
without the tail restriction.
\end{lemma}

\begin{proof}
For $t\ge1$, Cauchy--Schwarz gives the square bound
$\int_{|x|>R}w^2d\ga_n$ times the mass, which is its cost.
For $0\le t<1$ write $x=z+tu$ and
$m(z)=\int_0^1h(z+tu)\phi_1(t)dt$.
Every density $v\le\phi_0$ on the positive halfline with mass $m$
satisfies
\[
 \int tv(t)dt=\int_0^\infty\int_s^\infty v(t)dt\,ds
 \ge\int_0^{m/\phi_0}(m-\phi_0s)ds=\frac{m^2}{2\phi_0}.
\]
Thus
\begin{equation}\label{eq:bathtub}
 m(z)^2\le2\phi_0\int_0^1t h(z+tu)\phi_1(t)dt.
\end{equation}
In the tail, $|z|>\sqrt{R^2-1}$ and
$w(z+tu)\le2e^{1/16}(1+|z|^2)e^{|z|^2/16}$.
Apply Cauchy--Schwarz in $z$ and \eqref{eq:bathtub}.
The two remaining factors are the tails of $w^2$ against $\ga_n$ and
of $(1+|z|^2)^2e^{|z|^2/8}$ against $\ga_{n-1}$. Both vanish.
For $n=1$ the second tail is empty. Omitting tail restrictions gives
the global constant.
\end{proof}

\begin{lemma}\label{lem:kernel}
For $0<\rho\le1/16$, $L_\rho=(M_\rho-1)/\rho$ satisfies
\begin{equation}\label{eq:kernel-weight}
 |L_\rho(x,y)|\le Cw(x)w(y).
\end{equation}
It converges to $x\cdot y$ uniformly on compact products.
\end{lemma}

\begin{proof}
Use $L_\rho=\int_0^1\partial_rM_r|_{r=s\rho}ds$.
The derivative is a quadratic polynomial with bounded coefficients
times $M_r$, and \eqref{eq:mehler} gives
$M_r\le C\exp[(|x|^2+|y|^2)/34]$ on this range.
The weight $w(x)w(y)$ dominates their product. The derivative at zero
gives the compact limit.
\end{proof}

Let $\Pi_0v=\int v\,d\ga_n$ and put
$\widetilde Q_\rho(v)=\sum_i\ip{v_i}{(T_\rho-\Pi_0)v_i}_{L^2(\ga_n)}$.
For balanced $A$, quadratic expansion gives
\begin{equation}\label{eq:comparison}
 S_\rho(A)-S_\rho(T)=\widetilde Q_\rho(e)-2\LL_\rho(A).
\end{equation}

\begin{corollary}\label{cor:exterior}
For all partitions and $0<\rho\le1/16$,
\begin{equation}\label{eq:tv-tail}
 \left(\sum_i\int_{|x|>R}w|e_i|d\ga_n\right)^2
 \le C\beta(R)\LL_\rho(A)/\rho.
\end{equation}
The corresponding global bound has a finite constant. Also
$\widetilde Q_\rho(e)/\rho\le C(\sum_i\int w|e_i|d\ga_n)^2$.
\end{corollary}

\begin{proof}
Sum Lemma~\ref{lem:weighted} over the $k(k-1)$ ordered transfers and
use \eqref{eq:regret-cost}. The interaction bound is
Lemma~\ref{lem:kernel} integrated against $e_i\ga_n$ in both variables.
\end{proof}

\section{From measurable transfers to facet fluxes}\label{sec:compactness}

The following proposition is the nonlinear step. There is no regularity
assumption on the perturbed partitions.

\begin{proposition}\label{prop:compactness}
Let $\rho_m\downarrow0$ and let balanced partitions $A^{(m)}$ of
$\R^n$ approach $T$ in Gaussian label measure. Suppose
$s_m^2=\LL_{\rho_m}(A^{(m)})/\rho_m>0$ and $s_m\to0$.
Along a subsequence there is $f\in\HH_0$ such that
\begin{equation}\label{eq:flux-liminf}
 1\ge\frac\ell2\norm f_\HH^2,\qquad
 \frac{\widetilde Q_{\rho_m}(e^{(m)})}{\rho_ms_m^2}
       \longrightarrow\norm{\DD f}_F^2.
\end{equation}
If $B_m-B_*\perp\VV$ for every $m$, then $\DD f\perp\VV$.
\end{proposition}

\begin{proof}
Normalize the signed measures by $\mu_{m,i}=e_i^{(m)}\ga_n/s_m$.
Corollary~\ref{cor:exterior} bounds their weighted total variations and
gives the uniform tail bound
\begin{equation}\label{eq:normalized-tightness}
 \sup_m\sum_i\int_{|x|>R}w\,d|\mu_{m,i}|
       \le C\sqrt{\beta(R)}\longrightarrow0.
\end{equation}

\smallskip\noindent\emph{The limiting measures lie on facets.}
On each whole pair plane define
\begin{equation}\label{eq:theta}
 \theta_{m,ij}(z)=\frac1{s_m\phi_0}
                  \int_0^\infty h_{ij}^{(m)}(z+tn_{ij})\phi_1(t)dt.
\end{equation}
These nonnegative densities are bounded in $L^2(\nu_{ij})$.
For the normal integral over $t<1$ use \eqref{eq:bathtub}; the square
of its part over $t\ge1$ is bounded by a constant times that mass.
Their sum is controlled by $\int\min\{t,1\}h\phi_1dt$.
Equation~\eqref{eq:regret-cost} proves the claim after division by
$s_m^2$. Pass simultaneously to weak limits $\theta_{ij}\ge0$.

For every fixed $b>0$, \eqref{eq:regret-cost} gives
\begin{equation}\label{eq:far-normal}
 s_m^{-1}\int_{t\ge b}h_{ij}^{(m)}d\ga_n\le C_bs_m\to0.
\end{equation}
For a compact continuous test function, its value at $z+tn_{ij}$ can
therefore be replaced by its value at $z$: use uniform continuity for
$t<b$, bounded normalized mass, and then let $b\downarrow0$.
Thus the ordered measures converge to $\theta_{ij}\nu_{ij}$.
Their support is in $F_{ij}$, since any compact subset of
$H_{ij}\setminus F_{ij}$ has positive distance from $T_i$ and requires
a fixed normal displacement. Junctions are plane-null sets and acquire
no mass from these $L^2$ densities; there are no junctions when $n=1$.

Put $f_{ij}=\theta_{ji}-\theta_{ij}$ for $i<j$ and extend
antisymmetrically. Then
\begin{equation}\label{eq:measure-limit}
 \mu_{m,i}\rightharpoonup\mu_i
       =\sum_{j\ne i}f_{ij}\nu_{ij}|_{F_{ij}}.
\end{equation}
The weighted tails pass total masses and first moments. Balance gives
$f\in\HH_0$, and the asserted moment orthogonality passes to $\DD f$.

\smallskip\noindent\emph{The cost is at least half the squared flux.}
Fix a compact tangential cutoff $0\le\chi\le1$ and $0<b<1$.
Exact vanishing on the pair plane and local gradient convergence give
\[
 U_{\rho_m,ij}(z+tn_{ij})\ge\rho_m(\ell-o(1))t
 \qquad(0\le t\le b)
\]
on its compact normal cylinder. This is a multiplicative estimate, even
for transfers much thinner than $\rho_m$. Write
$m_{m,ij}^b(z)=\int_0^b h_{ij}^{(m)}(z+tn_{ij})\phi_1(t)dt$.
Discard nonnegative loss outside the cylinders and apply the density
bound to obtain explicitly
\[
 1\ge\frac{\ell-o(1)}2\sum_{i\ne j}
       \int\chi(z)\left(\frac{m_{m,ij}^b(z)}{s_m\phi_0}\right)^2d\nu_{ij}(z).
\]
The truncated densities have the same weak $L^2$ limits as
\eqref{eq:theta}: their omitted normalized mass tends to zero by
\eqref{eq:far-normal}, and their $L^2$ norms stay bounded.
Approximation of $L^2$ tests by bounded tests justifies this passage.
Weak lower semicontinuity and exhaustion by the cutoffs give
\begin{equation}\label{eq:ordered-liminf}
 1\ge\frac\ell2\sum_{i\ne j}\norm{\theta_{ij}}_{L^2(\nu_{ij})}^2
       \ge\frac\ell2\norm f_\HH^2.
\end{equation}
The last step uses $(\theta_{ji}-\theta_{ij})^2\le
\theta_{ij}^2+\theta_{ji}^2$. The factor is exactly $\ell/2$ because
$d\nu_{ij}=\phi_0d\ga_{n-1}$.

\smallskip\noindent\emph{Only the first-moment interaction survives.}
The normalized interaction can be written as
\[
 \sum_i\iint L_{\rho_m}(x,y)d\mu_{m,i}(x)d\mu_{m,i}(y).
\]
Insert continuous compact cutoffs in both variables. On their product,
$L_{\rho_m}\to x\cdot y$ uniformly, and the resulting finite sum of
products of first moments converges by \eqref{eq:measure-limit}.
The omitted terms are bounded by \eqref{eq:kernel-weight} and
\eqref{eq:normalized-tightness}, uniformly in $m$. Letting the cutoffs
exhaust space, and using $|x|\le w(x)$, gives
\[
 \sum_i\iint L_{\rho_m}(x,y)d\mu_{m,i}(x)d\mu_{m,i}(y)
   \longrightarrow\sum_i\left|\int x\,d\mu_i\right|^2
   =\norm{\DD f}_F^2.
\]
This proves \eqref{eq:flux-liminf}.
\end{proof}

\section{The contradiction and dimension reduction}

\begin{theorem}\label{thm:rigidity}
For each fixed $k\ge2$, there is $\rho_k>0$ such that every balanced
partition $A$ of $\R^{k-1}$ with $S_\rho(A)\ge S_\rho(T)$ and
$0<\rho\le\rho_k$ is a regular-simplex partition.
\end{theorem}

\begin{proof}
Otherwise choose $\rho_m\downarrow0$ and nonsimplex partitions
$A^{(m)}$ with $S_{\rho_m}(A^{(m)})\ge S_{\rho_m}(T)$.
The Hermite expansion and Parseval give
\[
 S_\rho(A)=1/k+\rho W_1(A)+R_\rho(A),\qquad
 0\le R_\rho(A)\le(1-1/k)\rho^2.
\]
Thus $0\le\delta_m=F-W_1(A^{(m)})\le(1-1/k)\rho_m$.
Choose the closest regular moment matrix and rotate it to $B_*$.
Proposition~\ref{prop:alignment} gives simultaneously
\begin{equation}\label{eq:gauge}
 \ga_n(I_{A^{(m)}}\ne I_T)\to0,\qquad B_m-B_*\perp\VV.
\end{equation}
Rotations preserve balance and noise stability.

The score loss is positive, since the aligned partition is not $T$.
Set $s_m^2=\LL_{\rho_m}(A^{(m)})/\rho_m$.
Equations~\eqref{eq:comparison} and Corollary~\ref{cor:exterior} give
\begin{equation}\label{eq:scale}
 s_m^2\le\frac{\widetilde Q_{\rho_m}(e^{(m)})}{2\rho_m}
 \le C\left(\sum_i\int w|e_i^{(m)}|d\ga_n\right)^2\longrightarrow0.
\end{equation}
Here \eqref{eq:gauge} and integrability of $w$ imply the last limit.
Proposition~\ref{prop:compactness} now supplies a flux with
\[
 \norm{\DD f}_F^2\ge2\ge\ell\norm f_\HH^2,
 \qquad f\in\HH_0,\quad\DD f\perp\VV.
\]
Hence $Q(f)\le0$. Lemma~\ref{lem:interface} forces $Q(f)=0$ and
$f\in\RR$, so $\DD f\in\VV$. Orthogonality then gives
$\DD f=0$, contradicting its squared norm being at least $2$.
\end{proof}

\begin{proof}[Proof of Theorem~\ref{thm:main}]
A maximizer exists at each fixed positive correlation with the prescribed
masses. For completeness, the fractional feasible set
\[
 \{(u_i)\in L^2(\ga_d)^k:0\le u_i\le1,\ \sum_i u_i=1,
                                  \ \int u_i d\ga_d=1/k\}
\]
is weakly compact, and $T_\rho$ is compact: its Hermite eigenspaces have
finite multiplicity and eigenvalues $\rho^m\to0$. The quadratic objective
is therefore weakly continuous. A fractional maximizer must be ordinary.
Otherwise two components are bounded away from both zero and one on a
set of positive measure. A nonzero bounded mean-zero $g$ on that set gives
feasible perturbations $u\pm\varepsilon g(e_i-e_j)$ whose average
objective is larger, since $\ip g{T_\rho g}>0$ by the positive Hermite
eigenvalues. This proves existence without dropping the mass constraints.

By Heilman--Tarter \cite[Theorem~1.9]{HT}, every maximizer in $d\ge n$
is cylindrical over $\R^n$ after rotation. In their notation the required
cell-count condition is $k\le d+1$, exactly our range. The reduced
partition has the same masses and stability, so
Theorem~\ref{thm:rigidity} makes it a simplex for $0<\rho\le\rho_k$.
Every equality case is itself a maximizer and undergoes the same reduction.
\end{proof}

\section{Why the first-moment deficit is a different scale}

The threshold in Theorem~\ref{thm:main} is qualitative, and the proof
does not compare partitions at every positive correlation. A separate
limitation explains the use of exact score loss rather than first-moment
deficit.

\begin{proposition}\label{prop:obstruction}
Fix $k\ge2$ and $0<\rho<1$. There is no $c(\rho)>0$ such that
\[
 S_\rho(T)-S_\rho(A)\ge c(\rho)\rho\,[F-W_1(A)]
\]
for all balanced partitions with positive first-moment deficit, even in
arbitrarily small Gaussian-measure neighborhoods of $T$ in $\R^{k-1}$.
\end{proposition}

\begin{proof}
Choose a unit direction $v$ strictly inside $T_i$ and a distinct label
$j$. Let $G$ reflect in $H_{ij}$. For large $R$, the ball
$E_R=B(Rv,1)$ lies in $T_i$, and $GE_R\subset T_j$ has the same
mass $p_R$. Exchange the two balls. This gives a balanced partition
$A_R$ approaching $T$. Put
$a_R=\int_{E_R}n_{ij}\cdot x\,d\ga_n$; then
$a_R\ge(Rn_{ij}\cdot v-1)p_R$ and $a_R\to0$.
The two moment changes are $-2a_Rn_{ij}$ and $2a_Rn_{ij}$, so
\[
 F-W_1(A_R)=4\ell a_R-8a_R^2\ge2\ell a_R>0
\]
for large $R$. Only points in the two balls change labels; agreement can
change only if at least one observation lies there. Hence
$|S_\rho(T)-S_\rho(A_R)|\le4p_R$. Dividing by the displayed
positive deficit times $\rho$ gives a ratio tending to zero.
\end{proof}

For these remote changes, linear score loss grows like distance times
mass, while the bounded OU score loss saturates. The normalization in
\eqref{eq:scale} keeps precisely this bounded-score cost in the
compactness argument. It imposes no lower bound on the size or location
of the perturbation.

\paragraph{AI assistance and review.}
ChatGPT (OpenAI) assisted with mathematical arguments, literature checks,
exposition, proof review and verification code.

\begin{thebibliography}{9}
\raggedright
\bibitem{CW}
C.~Christof and G.~Wachsmuth,
\emph{No-gap second-order conditions via a directional curvature functional},
SIAM J. Optim. \textbf{28} (2018), 2097--2130.
\href{https://doi.org/10.1137/17M1140418}{doi:10.1137/17M1140418}.
Theorem numbering refers to \href{https://arxiv.org/abs/1707.07579v2}{arXiv:1707.07579v2}.

\bibitem{Heilman14}
S.~Heilman, \emph{Euclidean partitions optimizing noise stability},
Electron. J. Probab. \textbf{19} (2014), no.~71, 1--37.
\href{https://doi.org/10.1214/EJP.v19-3083}{doi:10.1214/EJP.v19-3083}.

\bibitem{HT}
S.~Heilman and A.~Tarter,
\emph{Three candidate plurality is stablest for small correlations},
Forum Math. Sigma \textbf{9} (2021), e65.
\href{https://doi.org/10.1017/fms.2021.56}{doi:10.1017/fms.2021.56}.

\bibitem{MN}
E.~Milman and J.~Neeman,
\emph{The Gaussian double-bubble and multi-bubble conjectures},
Ann. of Math. (2) \textbf{195} (2022), 89--206.
\href{https://doi.org/10.4007/annals.2022.195.1.2}{doi:10.4007/annals.2022.195.1.2}.

\bibitem{Zhang}
Y.~Zhang (mxym/math research project),
\emph{Equal Gaussian cells: the sharp simplex first-moment theorem},
preprint, 8 October 2026.
Release \texttt{gaussian-balanced-simplex-all-k-v1},
commit \texttt{ebfd90558439afca21babeb7b77c8d8837ea745f}.
\url{https://github.com/mxym/math/releases/tag/gaussian-balanced-simplex-all-k-v1}.
\end{thebibliography}
\end{document}
